\subsection{Bart\'ok et al.\ (2010): Gaussian Approximation Potentials}
\label{sec:appendix-example-bartok2010}

\paragraph{Q1 -- Bibliographic identification.}
Bart\'ok, Payne, Kondor, and Cs\'anyi introduce the Gaussian
Approximation Potential (GAP): Gaussian-process regression on a
bispectrum-based descriptor of the local atomic environment.
The method is demonstrated on bulk semiconductors (C, Si, Ge), GaN,
and $\alpha$-Fe; we encode bulk silicon as the lead system.
Published in \emph{Phys.\ Rev.\ Lett.} 104, 136403 (2010);
arXiv:0910.1019.
\lstinputlisting[language=turtle]{artifacts/kg/listings/bartok2010/q01-bibliographic.ttl}

\paragraph{Q2 -- Material system.}
The lead demonstration system encoded here is bulk silicon in its
diamond-cubic phase. The paper also reports analogous GAP fits for
diamond (carbon), germanium, GaN, and $\alpha$-iron; those systems
are flagged in Q12.
\lstinputlisting[language=turtle]{artifacts/kg/listings/bartok2010/q02-material.ttl}

\paragraph{Q3 -- Reference calculation method.}
Reference data are produced from plane-wave DFT calculations.
\lstinputlisting[language=turtle]{artifacts/kg/listings/bartok2010/q03-reference-method-type.ttl}

\paragraph{Q4 -- Reference settings.}
The DFT engine is CASTEP. Specific exchange--correlation functional,
plane-wave cutoff, k-mesh, and pseudopotential type are reported
in the Supplementary Information rather than the main letter, so
they are not encoded here. \emph{Not reported in the main text.}
\lstinputlisting[language=turtle]{artifacts/kg/listings/bartok2010/q04-reference-settings.ttl}

\paragraph{Q5 -- Training dataset.}
Reference configurations are produced by random atomic and lattice
displacements (up to 0.2\,\AA) of equilibrium 2-, 8-, 16-, and
64-atom cubic Si unit cells, with energies and forces from CASTEP.
The model is later extended with random displacements around a
diamond vacancy and a graphite-to-diamond transition path.
$\dlRole{numConfigurations}$ is not enumerated in the main text;
provenance is \emph{in-house}.
\lstinputlisting[language=turtle]{artifacts/kg/listings/bartok2010/q05-dataset.ttl}

\paragraph{Q6 -- Sampling strategies.}
A single sampling strategy applies: random atomic and lattice
displacements of small cubic unit cells, supplemented at the end of
the paper with isolated-defect (vacancy) and reaction-pathway
(graphite$\to$diamond) configurations.
\lstinputlisting[language=turtle]{artifacts/kg/listings/bartok2010/q06-sampling.ttl}

\paragraph{Q7 -- MLIP method, components, implementation.}
GAP is Gaussian-process regression on a four-dimensional
bispectrum descriptor of the atomic neighbour density inside a
smooth cosine cutoff; total energy is the sum over atoms.
Training is sparse GP regression with $M$ randomly selected
representative atomic-neighbourhood configurations. The
implementation is the QUIP/GAP suite (\texttt{libatoms.org}).
\lstinputlisting[language=turtle]{artifacts/kg/listings/bartok2010/q07-method.ttl}

\paragraph{Q8 -- Hyperparameter settings.}
The main letter does not tabulate explicit numerical values for the
descriptor cutoff $r_c$, bispectrum truncation $J_{\max}$, or
sparse-point count $M$ (these are reported in the SI). We
therefore record no $\HyperparameterSettingC$ instances in the
canonical encoding and flag them as Q12 gaps.
\lstinputlisting[language=turtle]{artifacts/kg/listings/bartok2010/q08-hyperparameter-settings.ttl}

\paragraph{Q9 -- MLIP run and trained model.}
A single sparse-GP training run produces one GAP for diamond-cubic
Si. We attach the per-atom inference time
($\dlRole{inferenceTimePerAtom}=0.01$\,s/atom/timestep, single CPU
core) to the trained model.
\lstinputlisting[language=turtle]{artifacts/kg/listings/bartok2010/q09-run-and-model.ttl}

\paragraph{Q10 -- Benchmark results.}
The headline accuracy figure is an energy RMSE bound of
$<1$\,meV/atom on near-bulk Si configurations, with the same bound
holding for diamond and iron. We record the bound as the value
($1$\,meV/atom) on the silicon $\BenchmarkResultC$. Phonon
dispersions, elastic constants, and defect-formation energies are
also reported but in derived form rather than as separate
$\BenchmarkResultC$ instances.
\lstinputlisting[language=turtle]{artifacts/kg/listings/bartok2010/q10-benchmark-results.ttl}

\paragraph{Q11 -- Computational resources.}
The paper reports a per-atom inference cost of 0.01\,s/atom/timestep
on a single CPU core (encoded on the trained model in Q9). It also
notes that the 216-atom unit-cell timestep takes 191\,s/atom in
CASTEP --- about 20{,}000$\times$ slower than GAP. Training duration,
training hardware, peak memory, and inference hardware are not
reported.
\lstinputlisting[language=turtle]{artifacts/kg/listings/bartok2010/q11-resources.ttl}

\paragraph{Q12 -- Gaps.}
Beyond Q11, the main letter does not tabulate: (i) DFT-side
parameters (exchange--correlation functional, plane-wave cutoff,
k-mesh, pseudopotential type) which live in the SI; (ii) the
explicit number of training configurations
$\dlRole{numConfigurations}$; (iii) explicit
$\HyperparameterSettingC$ values for $r_c$, $J_{\max}$, $M$,
GP kernel parameters, and per-channel noise; (iv) the GP
hyperparameter optimiser; (v) any $\dlRole{frozenCore}$ treatment
beyond the CASTEP pseudopotential. The companion systems (C, Ge,
GaN, $\alpha$-Fe) are not encoded here; each would live in a
sibling file under the same $\BenchmarkStudyC$ aegis.
